\documentclass[12pt]{article}
\usepackage{amsmath, amssymb, amsthm}
\usepackage{graphicx}
\title{Discrete Math HW4}
\author{Anatolii Larkin}

\begin{document}

\maketitle


\begin{enumerate}

\item Suppose you have a deck of cards, each card has a number on one side, and
the opposite side is colored red or green. You want to check the following rule:
«if a card has an even number on one side, then the opposite side is colored
red».
You see 4 cards on the desk, they show: number 6, number 1, red, green.
What is the minimum number of cards (and which ones) you have to turn
over to look at to check if the rule holds true?

Check card number 6. If it's red, the rule holds so far. If it is green, the rule doesn't hold.

Check green card. If the numbers is odd, the rule holds. If it is even, the rule doesn't hold.

No more cards are required to check if this rule holds.

Card number 1 can be any color, as it would not violate the rule.

Red card can have any number, as it would not violate the rule.




\item Build the truth table for a boolean function \(f(a,b,c) = ((b \land a) \oplus a) \lor (\neg a) \lor (c \land a)\)

$(f(a,b,c) = ((b \land a) \oplus a) \lor (\neg a) \lor (c \land a)\) = (\neg b \land a) \lor (\neg a) \lor (c \land a) = (\neg b \lor \neg a) \lor (c \land a) = \neg (b \land a) \lor (c \land a)$

\begin{tabular}{c|c|c|c}
\(a\) & \(b\) & \(c\) & \(f(a, b, c)\) \\
\hline
T & T & T & T \\
T & T & F & F \\
T & F & T & T \\
T & F & F & T \\
F & T & T & T \\
F & T & F & T \\
F & F & T & T \\
F & F & F & T \\
\end{tabular}


\item Determine, which of the functions will be constant 1 (always return 1), and which will be constant 0.

\begin{enumerate}
    \item \((x \implies (y \implies z)) \implies ((x \implies y) \implies z)\)

    Since $a \implies b \equiv \neg a \lor b$:

    $(x \implies (y \implies z)) \implies ((x \implies y) \implies z) \equiv (\neg x \lor \neg y \lor z) \implies ((x \land \neg y) \lor z) \equiv \neg(\neg x \lor \neg y \lor z) \lor ((x \land \neg y) \lor z) \equiv (x \land y \land \neg z) \lor ((x \land \neg y) \lor z) = z \lor (x \land y) \lor (x \land \neg y) = z \lor (x \land (y \lor \neg y)) = z \lor x \in \{0,1\}$ (not constant)
    
    \item $((x \implies y) \land (y \implies z)) \implies (x \implies z)
    \equiv \neg ((\neg x \lor y) \land (\neg y \lor z)) \lor (\neg x \lor z)
    \equiv \neg (\neg x \lor y) \lor \neg (\neg y \lor z) \lor \neg x \lor z
    \equiv (x \land \neg y) \lor (y \land \neg z) \lor \neg x \lor z
    \equiv (x \land \neg y) \lor \neg x \lor (y \land \neg z) \lor z
    \equiv ((x \lor \neg x) \land (\neg y \lor \neg x)) \lor ((y \lor z) \land (\neg z \lor z))
    \equiv (1 \land (\neg y \lor \neg x)) \lor ((y \lor z) \land 1)
    \equiv \neg x \lor \neg y \lor y \lor z
    \equiv \neg x \lor 1 \lor z
    \equiv 1$
    
    \item \(x \land (x \implies y) \land (x \implies \neg y)\)

    $x \land (x \implies y) \land (x \implies \neg y) \equiv x \land (\neg x \lor y) \land (\neg x \lor \neg y) \equiv (x \land \neg x) \lor (x \land y) \lor (\not x \lor \not y) \equiv 0 \lor (x \land y) \land (\neg x \lor \neg y) \equiv 0 \lor (x \land y) \land \neg (x \land y) = 0 \lor 0 \equiv 0$
\end{enumerate}

\item Simplify the functions, using in the answer only operations \(\land\), \(\lor\) и \(\neg\) (conjunction, disjunction and negation).

(a) \((x \implies y) \implies y\)

Since $a \implies b \equiv \neg a \lor b$:

$(x \implies y) \implies y \equiv (\neg x \lor y) \implies y \equiv \neg (\neg x \lor y) \lor y \equiv (x \land \neg y) \lor y \equiv (y \lor \neg y) \land (x \lor y) \equiv T \land (x \lor y) \equiv x \lor y$


(b) \(\neg(x \implies \neg y)\)

$\neg(x \implies \neg y) \equiv \neg (\neg x \lor \neg y) \equiv x \land y$


\item Express this boolean function (defined by the following truth table), using \(\land\), \(\lor\) and \(\neg\). And simplify the answer.

\[
\begin{array}{c|c|c|c}
x & y & z & f(x, y, z) \\
\hline
0 & 0 & 0 & 0 \\
0 & 0 & 1 & 1 \\
0 & 1 & 0 & 0 \\
0 & 1 & 1 & 0 \\
1 & 0 & 0 & 0 \\
1 & 0 & 1 & 1 \\
1 & 1 & 0 & 0 \\
1 & 1 & 1 & 1
\end{array}
\]


$f(x,y,z) = 1 \equiv (\neg x \land \neg y \land z) \lor (x \land \neg y \land z) \lor (x \land y \land z) = (\neg y \land z) \land (\neg x \lor x) \lor (x \land y \land z) = (\neg y \land z) \lor (x \land y \land z) = z \land ((\neg y \lor y) \land (\neg y \lor z)) = z \land (\neg y \lor (x \land y)) = z \land ((\neg y \lor y) \land (x \lor \neg y)) = z \land (x \lor \neg y)$

\[
\begin{array}{c|c|c|c|c}
x & y & z & f(x, y, z) & $z \land (x \lor \neg y)$ \\
\hline
0 & 0 & 0 & 0 & 0\\
0 & 0 & 1 & 1 & 1\\
0 & 1 & 0 & 0 & 0\\
0 & 1 & 1 & 0 & 0\\
1 & 0 & 0 & 0 & 0\\
1 & 0 & 1 & 1 & 1\\
1 & 1 & 0 & 0 & 0\\
1 & 1 & 1 & 1 & 1
\end{array}
\]

\item Fill the blank spaces with universal and existence quantifiers to get true statements. Write the most general statements (i.e. choose universal quantifier if both solutions are possible)

\begin{itemize}
    \item \(\dots x \in \mathbb{N} : x > 10\)

    $\exists x \in \mathbb{N}: x > 10$
    
    \item \(\dots x \in \mathbb{N} \dots y \in \mathbb{N} : x + y = 10\)

    $\exists x \in \mathbb{N} \exists y \in \mathbb{N}: x + y = 10$
    \item \(\dots x \in \mathbb{N} \dots y \geq 1 \in \mathbb{N} : x \bmod y = 0\)

    $\forall x \in \mathbb{N} \exists y \geq 1 \in \mathbb{N}: x \mod y = 0$
    
    \item \(\dots x \in \mathbb{R} \dots y \in \mathbb{R} : x^2 + y^2 \geq 0\)

    $\forall x \in \mathbb{R} \forall y \in \mathbb{R}: x^2 + y^2 \geq 0$
    
    \item \(\dots x \in \mathbb{N} \dots y \in \mathbb{N} : x = \frac{(y + 77) \cdot 2 + 324}{2} - y\)

    $x = \frac{(y+77)*2+324}{2} - y \equiv \frac{2y + 478}{2}-y = x \equiv x = 239$

    $\exists x \in \mathbb{N} \forall y \in \mathbb{N}: x = 239$
    
\end{itemize}

\end{enumerate}
\end{document}